\chapter{Phases 0--1: the ODE prototype, identification, decidability}\label{ch:cII1}
\chapterbrief
{the toolkit of Part~I, Chapter~\ref{ch:cI1} (A1--A9). For the companion matrix: $P\in\C[\xi][\tau]$ monic in $\tau$ of degree $k$, $C(\xi)$ its companion matrix (ones on the superdiagonal, last row $(-p_0,\dots,-p_{k-1})$), $v(\tau)=(1,\tau,\dots,\tau^{k-1})^{\mathsf T}$.}
{to prove the remaining results of Chapters~3--5 of the paper: the one-variable prototype, the projectors of the companion matrix, the refinements of identification (evolutions, waves, symmetries, local data) and the decidability results beyond A8--A9.}
{Dim.~B1--B21. Two elementary facts are used throughout: (F0) $C(\xi)v(\tau)=\tau v(\tau)-P(\tau,\xi)e_k$, so $v(\tau_0)$ is an eigenvector when $P(\tau_0,\xi_0)=0$; (F1$'$) $e_k$ is a cyclic vector of $C$, so the minimal polynomial of $C$ is $P$.}

%=====================================================================
\section{The one-variable prototype and the companion matrix}

\begin{dimcard}{B1}{R4-T1}{ODE identification}{\proved}{prop:ode-ident}
Let $u\in\ker p(D)$, $D=d/dx$, $\deg p=m$. The minimal annihilator of $u$ has degree $\dimc\spn\{u,u',u'',\dots\}$, the number of basis functions $x^je^{\lambda x}$ actually present in $u$. It equals $p$ iff the Wronskian of $u,\dots,u^{(m-1)}$ does not vanish.
\end{dimcard}
\begin{proof}
\step{1}{The derivative space}
In one variable every ideal of $\C[\xi]$ is principal, so $\Ann(u)=(q)$ with $q$ the minimal annihilator, and $\C[\xi]u\cong\C[\xi]/(q)$ has dimension $\deg q$.

\step{2}{Which basis functions are present}
Write $u=\sum_jq_j(x)e^{\lambda_jx}$. By A2(c) in one variable, $(D-\lambda_j)^{\deg q_j+1}$ kills the $j$-th summand and no lower power does, so $q=\prod_j(\xi-\lambda_j)^{\deg q_j+1}$, of degree equal to the number of functions $x^le^{\lambda_jx}$ present.

\step{3}{The Wronskian}
$q=p$ iff $\deg q=m$ iff $u,\dots,u^{(m-1)}$ are linearly independent iff their Wronskian is not identically zero.
\whyworks{in one variable the derivative space of $u$ is the quotient by its minimal annihilator.}
\end{proof}
\uses{A2.}

%---------------------------------------------------------------------
\begin{dimcard}{B2}{II-L1}{First-order reduction as a module equivalence}{\pw}{}
(a) $\ring^{1\times k}/\ring^{1\times k}(\tau I-C(\xi))\cong\ring/(P)$, with $\bar e_1\mapsto\bar1$. (b) The $t$-jet map $J:u\mapsto(u,\partial_tu,\dots,\partial_t^{k-1}u)$ is a bijection from $\ker P$ onto the solutions of $(\partial_t-C(\partial_x))U=0$, with inverse $U\mapsto U_1$, and $J(f(\partial)w)=f(\partial_x,C(\partial_x))J(w)$.
\end{dimcard}
\begin{proof}
\step{1}{(a) The presentation}
Map $e_i\mapsto\tau^{i-1}$. The rows $\tau e_i-e_{i+1}$ go to $0$, and the last row $\tau e_k+\sum_jp_je_{j+1}$ goes to $P\equiv0$, so the map passes to the quotient and is onto. In the quotient $e_{i+1}\equiv\tau e_i$, so the module is cyclic, generated by $\bar e_1$, and its annihilator is exactly $(P)$.

\step{2}{(b) The jet map}
If $Pu=0$, $U=J(u)$ satisfies $\partial_tU_i=U_{i+1}$ and $\partial_tU_k=-\sum_jp_j(\partial_x)U_{j+1}$: this is $(\partial_t-C(\partial_x))U=0$. Conversely the first $k-1$ equations give $U_i=\partial_t^{i-1}U_1$ and the last gives $P(\partial)U_1=0$. Intertwining holds for $f=\tau$ by the equation itself, and for $f=\xi$ trivially.
\whyworks{the companion system is a presentation of the same cyclic module $\ring/(P)$.}
\end{proof}
\uses{none.}

%---------------------------------------------------------------------
\begin{dimcard}{B3}{II-T1}{When the companion form exists}{\pw}{}
(i) $\ring/(P)$ is a finitely generated $\C[\xi]$-module iff the leading coefficient of $P$ in $\tau$ is a nonzero constant; then it is free with basis $1,\tau,\dots,\tau^{k-1}$ and $\tau$ acts by $C(\xi)^{\mathsf T}$. (ii) For $P,Q$ monic in $\tau$: $C_P\sim C_Q\iff P=Q\iff\ring/(P)\cong\ring/(Q)$. (iii) If the principal symbol does not vanish at $(1,0)$, (i) holds with $k=\deg P$; the converse is false (heat).
\end{dimcard}
\begin{proof}
\step{1}{(i)}
For monic $P$, division by $P$ in $\C[\xi][\tau]$ gives the basis. Conversely, if the module is generated by $g_1,\dots,g_n$ over $\C[\xi]$, then $\tau g_i=\sum a_{ij}g_j$, and the determinant trick gives a monic $h(\tau)=\det(\tau I-(a_{ij}))$ killing the module, so $h=Pq$ and $1=\mathrm{lc}_\tau(P)\,\mathrm{lc}_\tau(q)$.

\step{2}{(ii)}
$\det(\tau I-C_P)=P$, similarity preserves it, and an isomorphism of modules preserves the annihilator $(P)$.

\step{3}{(iii)}
The coefficient of $\tau^m$ is $\sigma_m(1,0)$. Counterexample to the converse: $P=\tau-\xi^2$ has $\sigma_2(1,0)=0$ but a companion form of size one.
\whyworks{a companion form is a choice of a non-characteristic time direction.}
\end{proof}
\uses{B2.}

%---------------------------------------------------------------------
\begin{dimcard}{B4}{II-P2}{The principal symbol sees only infinity}{\pw}{}
(a) $\sigma_m(P+Q)=\sigma_m(P)$ if $\deg Q<m$: the principal symbol is not a complete invariant. (b) The projective closure of $\Var(P)$ meets the line at infinity exactly at the zeros of $\sigma_m$. (c) $\sigma_m$ transforms as a binary form under linear changes of variables.
\end{dimcard}
\begin{proof}
\step{1}{(a)}
The top-degree part does not change. Example: $\partial_t-\partial_x^2$ and $\partial_t-\partial_x^2+c$ have the same $\sigma_2=-\xi^2$ but different kernels.
\step{2}{(b), (c)}
$P^h(\tau,\xi,\zeta)=\sum_d\zeta^{m-d}\sigma_d$, so at $\zeta=0$ only $\sigma_m$ remains; the top-degree part commutes with linear substitution.
\whyworks{the principal symbol is the restriction of the equation to the line at infinity.}
\end{proof}
\uses{none.}

%---------------------------------------------------------------------
\begin{dimcard}{B5}{II-T2}{Projectors of the companion matrix, and their poles}{\proved}{prop:proj}
Let $P=Q_1\cdots Q_s$ with $Q_j$ monic in $\tau$ and pairwise coprime over $K=\C(\xi)$; $B_j=P/Q_j$, $r_j=\mathrm{Res}_\tau(Q_j,B_j)\in\C[\xi]$.
\begin{enumerate}[label=(\alph*)]
\item $r_j\ne0$, and $U_jQ_j+V_jB_j=r_j$ with $U_j,V_j\in\C[\xi][\tau]$.
\item $\Pi_j=(V_jB_j)(\xi,C(\xi))/r_j(\xi)$ are orthogonal idempotents commuting with $C$, summing to $I$; $\Pi_j$ projects onto $\ker Q_j(C)$ along $\ker B_j(C)$. At a point $\xi_0$ with $r_j(\xi_0)\ne0$ the same holds with kernel $\ker B_j(\xi_0,C_0)$.
\item The poles of $\Pi_j$ are exactly the $\xi$-coordinates of the points of $\Var(Q_j)\cap\Var(B_j)$.
\item On solutions: if $u=\sum_iu_i$ with $u_i\in\ker Q_i$, then $(V_jB_j)(\partial)u=r_j(\partial_x)u_j$.
\end{enumerate}
\end{dimcard}
\begin{proofidea}
This is the Chinese remainder theorem over the field $\C(\xi)$, evaluated at the companion matrix; the resultant is the common denominator, and it vanishes exactly where two factor curves meet.
\end{proofidea}
\begin{proof}
\step{1}{(a) Resultant identity}
$Q_j$ and $B_j$ are monic and coprime over $K$, so their resultant is nonzero, and Cramer's rule on the Sylvester matrix gives $U_jQ_j+V_jB_j=r_j$ with coefficients in $\C[\xi]$.

\step{2}{(b) Idempotents}
$e_j=V_jB_j/r_j\equiv1\pmod{Q_j}$ and $\equiv0\pmod{Q_i}$ for $i\ne j$. So $e_j^2\equiv e_j$, $e_ie_j\equiv0$, $\sum e_j\equiv1$ modulo $P$. Since $P(C)=0$ (F1$'$, Cayley--Hamilton), $f\mapsto f(C)$ is a ring homomorphism from $K[\tau]/(P)$, and the $\Pi_j=e_j(C)$ inherit these relations. If $Q_j(C)w=0$ then $\Pi_jw=w$; and $Q_j(C)\Pi_j=(V_jP)(C)/r_j=0$. So the image is $\ker Q_j(C)$; $B_j(C)(I-\Pi_j)=(U_jP)(C)/r_j=0$ gives the kernel.

\step{3}{(c) Poles}
If no point of $\Var(Q_j)\cap\Var(B_j)$ has abscissa $\xi_0$, the specialised polynomials are coprime and $r_j(\xi_0)\ne0$. Conversely, if $\Pi_j$ were holomorphic at a common point $(\tau_0,\xi_0)$, its value $\Pi_0$ would be an idempotent commuting with the cyclic $C_0$, hence a polynomial in $C_0$, hence a sum of spectral projectors; but $\tau_0$ is a root of both $Q_j$ and $B_j$ at $\xi_0$, and $C_0$ has a single Jordan block there, so $\Pi_0$ cannot both keep and kill it.

\step{4}{(d) On solutions}
For $i\ne j$, $(V_jB_j)(\partial)u_i=0$ because $Q_i\mid B_j$. And $(V_jB_j)(\partial)u_j=(r_j-U_jQ_j)(\partial)u_j=r_j(\partial_x)u_j$.
\whyworks{the resultant is the universal denominator of the Chinese remainder idempotents.}
\end{proof}
\uses{B2.}

\begin{example}{Wave}
For $(\tau-\xi)(\tau+\xi)$: $r=2\xi$ and
\[
\Pi_1=\begin{pmatrix}\tfrac12&\tfrac1{2\xi}\\[2pt]\tfrac\xi2&\tfrac12\end{pmatrix},
\]
with a pole at $\xi=0$, the abscissa of the crossing point of the two lines.
\end{example}

%---------------------------------------------------------------------
\begin{dimcard}{B6}{II-T3}{Polynomial projectors iff disjoint varieties}{\proved}{prop:poly-proj}
\phantomsection\label{dim:II-C2}%
For $P=P_1P_2$ and any $\Omega$, equivalent: (i) $\Var(P_1)\cap\Var(P_2)=\emptyset$; (ii) $(P_1,P_2)=\ring$; (iii) $\ring/(P)\cong\ring/(P_1)\times\ring/(P_2)$; (iv) some $E_1\in\ring$ is the identity on $\ker P_1$ and zero on $\ker P_2$; (v) $\ker P_1\cap\ker P_2=\{0\}$. Under (ii), $\ker P=\ker P_1\oplus\ker P_2$ on every $\Omega$. \textbf{(II-C2)} If (i) holds for nonconstant $P_1,P_2$, their principal symbols share a linear factor.
\end{dimcard}
\begin{proof}
\step{1}{(i)$\iff$(ii)$\iff$(iii)}
Weak Nullstellensatz; Chinese remainder theorem for comaximal ideals.

\step{2}{(ii)$\Rightarrow$(iv),(v), and the splitting}
From $aP_1+bP_2=1$ put $E_1=bP_2$: $E_1u_1=(1-aP_1)u_1=u_1$ and $E_1u_2=0$. Every $u\in\ker P$ splits as $(bP_2)(\partial)u+(aP_1)(\partial)u$, with the two pieces in $\ker P_1$ and $\ker P_2$.

\step{3}{(iv),(v)$\Rightarrow$(i)}
A common zero $(\tau_0,\xi_0)$ gives $w=e^{\tau_0t+\xi_0x}\in\ker P_1\cap\ker P_2$, $w\ne0$, contradicting (v) and, via $w=E_1w=0$, (iv).

\step{4}{(II-C2)}
Two projective plane curves always meet (B\'ezout). Under (i) they meet only at infinity, that is (B4(b)) at a common zero of the principal symbols.
\whyworks{a polynomial projector is a B\'ezout identity, which exists exactly when the curves never meet.}
\end{proof}
\uses{B4.}

%---------------------------------------------------------------------
\begin{dimcard}{B7}{II-T4}{Only factors give single-valued projectors}{\proved}{}
For $P$ monic in $\tau$ and a set $S$ of roots $\tau_a(\xi)$ (with full multiplicities), the spectral projector $\Pi_S$ of $C(\xi)$ is rational in $\xi$ iff $Q_S=\prod_{a\in S}(\tau-\tau_a)^{\mu_a}$ lies in $\C[\xi][\tau]$, that is, iff $S$ is a union of irreducible factors of $P$. There are exactly $2^r$ rational projectors, $r$ the number of distinct irreducible factors.
\end{dimcard}
\begin{proof}
\step{1}{Rational projector $\Rightarrow$ rational factor}
$\Pi_S=g(C)$ with $g$ of degree $<k$; since $I,\dots,C^{k-1}$ are independent with coefficients in $K$, Cramer gives $g\in K[\tau]$. Then $Q_S=\gcd(P,g-1)$ (with multiplicities) is computed by Euclid in $K[\tau]$, so $Q_S\in K[\tau]$, and by Gauss in $\C[\xi][\tau]$.

\step{2}{Rational factor $\Rightarrow$ rational projector}
B5(b) with $s=2$.
\whyworks{the branches of an irreducible factor are permuted by monodromy, so no single-valued projector separates them.}
\end{proof}
\uses{B5.}

%---------------------------------------------------------------------
\begin{dimcard}{B8}{II-P5}{Normality of the companion matrix}{\proved}{}
For a companion matrix $C$ of $p(s)=s^k+\sum_{j<k}p_js^j$, equivalent: (1) $C$ normal; (2) $C$ has an orthonormal eigenbasis; (3) the roots are simple and $\lambda_a\bar\lambda_b$ is a nontrivial $k$-th root of unity for all $a\ne b$; (4) $k=1$; or $k=2$ with $|p_0|=1$ and $\bar p_0p_1=-\bar p_1$; or $k\ge3$ with $p=s^k-c$, $|c|=1$. \textbf{The earlier claim ``normal iff $p=s^k-c$'' is false for $k=2$} \falsest: $s^2+s-1$ has the symmetric companion matrix $\begin{pmatrix}0&1\\1&-1\end{pmatrix}$.
\end{dimcard}
\begin{proof}
\step{1}{Normal forces simple roots}
A normal matrix is diagonalisable, and a cyclic one (F1$'$) is diagonalisable only with simple roots.

\step{2}{Orthogonality of eigenvectors}
With simple roots the eigenvectors are $v(\lambda_a)$, and
\[
\langle v(\lambda_a),v(\lambda_b)\rangle=\sum_{j<k}z^j,\qquad z=\lambda_a\bar\lambda_b,
\]
which vanishes iff $z\ne1$ and $z^k=1$. This proves (1)$\iff$(3).

\step{3}{$k\ge3$}
With three indices, $|\lambda_a|^2=|\lambda_a\lambda_b||\lambda_a\lambda_c|/|\lambda_b\lambda_c|=1$, so the roots are $\lambda_1$ times all $k$-th roots of unity and $p=s^k-\lambda_1^k$.

\step{4}{$k=2$}
Direct computation:
\[
CC^*-C^*C=\begin{pmatrix}1-|p_0|^2&-p_1\bar p_0-\bar p_1\\-p_0\bar p_1-p_1&|p_0|^2-1\end{pmatrix}.
\]
\whyworks{the eigenvectors of a companion matrix are Vandermonde vectors, orthogonal only for roots on a rotated circle of roots of unity.}
\end{proof}
\uses{none.}

%=====================================================================
\section{Refinements of identification}

\begin{dimcard}{B9}{II-C1}{The $t$-jet is a cyclic vector}{\proved}{prop:cyclic}
For $P$ monic in $\tau$ of degree $k$: $\Ann(u)=(P)$ iff no nonzero $(b_0,\dots,b_{k-1})\in\C[\xi]^k$ has $\sum_jb_j(\partial_x)\partial_t^ju=0$.
\end{dimcard}
\begin{proof}
\step{1}{Division by $P$}
This is A3(d): every $Q$ is $AP+\sum_{j<k}b_j(\xi)\tau^j$ uniquely, $Q\in\Ann(u)$ iff $\sum b_j(\partial_x)\partial_t^ju=0$, and $Q\in(P)$ iff all $b_j=0$.
\whyworks{modulo $P$, every operator is a $\C[\xi]$-combination of the first $k$ time derivatives.}
\end{proof}
\uses{A3.}

%---------------------------------------------------------------------
\begin{dimcard}{B10}{C3-T4}{Evolutions: Hermite interpolation}{\proved}{thm:evol}
Let $P=\tau-a(\xi)$. (a) $\Ann(u)=(P)\iff u\notin\EP$. (b) If $u=\sum_kq_ke^{\langle z_k,\cdot\rangle}$ (distinct $z_k$, $q_k\ne0$), $\delta_k=\deg_xq_k$, $D=\sum_k(\delta_k+1)$ and $\deg a\le n$, then
\[
\{b\in\C[\xi]_{\le n}:\ \tau-b\in\Ann(u)\}=a+\prod_k(\xi-\xi_k)^{\delta_k+1}\cdot\C[\xi]_{\le n-D}.
\]
\end{dimcard}
\begin{proof}
\step{1}{(a)}
A3(a) with one factor and $P/P_1=1$.
\step{2}{(b) Reduction to an ODE in $x$}
$\tau-b=P+(a-b)(\xi)$, so $\tau-b\in\Ann(u)$ iff $c(\partial_x)u=0$ with $c=a-b$, $\deg c\le n$.
\step{3}{(b) Triangular conditions}
By A2(a),(b), $c(\partial_x)u=0$ iff for each $k$
\[
c(\xi_k+\partial_x)q_k=\sum_j\frac{c^{(j)}(\xi_k)}{j!}\,\partial_x^jq_k=0 .
\]
Since $\partial_x^jq_k$ has $x$-degree exactly $\delta_k-j$, the system is triangular and equivalent to $c^{(j)}(\xi_k)=0$ for $j\le\delta_k$, that is, $\prod_k(\xi-\xi_k)^{\delta_k+1}\mid c$.
\whyworks{each mode with a polynomial prefactor of degree $\delta$ fixes $\delta+1$ Hermite data of the symbol.}
\end{proof}
\uses{A2, A3.}

%---------------------------------------------------------------------
\begin{dimcard}{B11}{C3-T5}{Wave, Klein--Gordon, telegraph, beam}{\proved}{ex:wave-ident}
(a) Wave $P=(\tau-c\xi)(\tau+c\xi)$, $c\in\R\setminus\{0\}$, $u=f(x-ct)+g(x+ct)$: $\Ann(u)=(P)\iff f\notin\EP$ and $g\notin\EP$. (b) $\tau^2-\xi^2+\mu^2$ ($\mu\ne0$) and $\tau^2+\alpha\tau-c^2\xi^2$ ($\alpha,c\ne0$) are irreducible: identified iff $u\notin\EP$. (c) Beam $\tau^2+\xi^4$: identified iff $(\partial_t\pm i\partial_x^2)u\notin\EP$ for both signs.
\end{dimcard}
\begin{proof}
\step{1}{(a) The partial quotients}
\[
w_1=(\partial_t+c\partial_x)u=2c\,g'(x+ct),\qquad w_2=(\partial_t-c\partial_x)u=-2c\,f'(x-ct).
\]
\step{2}{(a) Pull-backs}
$F(x-ct)\in\EP(\Omega)$ iff $F\in\EP(J_F)$: if an exponential-polynomial $E$ equals $F(x-ct)$ on $\Omega$, then $(\partial_t+c\partial_x)E=0$ on $\Omega$, hence everywhere (A2(b)), so $E=F_0(x-ct)$ with $F_0\in\EP$, and $F=F_0$ on $J_F$. In one variable $f\in\EP\iff f'\in\EP$. Apply A3(a).
\step{3}{(b), (c)}
A nondegenerate conic is irreducible (bordered determinant $-\mu^2$, resp.\ $c^2\alpha^2/4$, nonzero). For the beam apply A3(a) to the two factors.
\whyworks{each factor of the wave operator carries one profile, and each profile must be non-exponential.}
\end{proof}
\uses{A2, A3.}

%---------------------------------------------------------------------
\begin{dimcard}{B12}{C3-C1}{The spectrum of an expansion is not intrinsic}{\falsest{} (counterexample proved)}{cex:C3C1}
The claim ``if $u=\sum_{z\in Z}c_ze^{\langle z,\cdot\rangle}$ on a bounded $\Omega$ with all $c_z\ne0$ and $Z\cap\Var(P_i)$ infinite for each $i$, then $\Ann(u)=(P)$'' is false.
\end{dimcard}
\begin{proof}
\step{1}{The counterexample}
$P=\tau+\xi$, $\Omega\subset\{|x-t|<\pi\}$. On $(-\pi,\pi)$,
\[
e^s=\sum_{k\in\Z}c_ke^{iks},\qquad c_k=\frac{\sinh(\pi)\,(-1)^k}{\pi(1-ik)}\ne0,
\]
in $L^2$, hence in $\mathcal D'(\Omega)$. So $u=e^{x-t}$ has infinitely many exponents $(-ik,ik)$ on $\tau+\xi=0$, but $\tau+1,\ \xi-1\in\Ann(u)$.
\step{2}{Diagnosis}
The expansion is not unique: $(-1,1)$ also lies on the line, giving a nontrivial relation. On a window of length $\ge2\pi$, Fourier uniqueness restores the criterion \led{C3-P4}.
\whyworks{on a bounded domain a function has many exponential expansions, so their exponents are not invariants.}
\end{proof}
\uses{none.}

%---------------------------------------------------------------------
\begin{dimcard}{B13}{C3-C2}{Exponential-polynomial data never identify}{\pw}{}
If $u\in\EP(\Omega)$ and $P(\partial)u=0$ with $P$ nonconstant, then $\Ann(u)\ne(P)$.
\end{dimcard}
\begin{proof}
\step{1}{Finite versus infinite}
$\Var(\Ann u)$ is finite (A2(c)), while $\Var(P)$ is infinite (A1(F)).
\whyworks{finitely many exponentials sit on finitely many points of a curve.}
\end{proof}
\uses{A1, A2.}

%---------------------------------------------------------------------
\begin{dimcard}{B14}{C3-P3}{Translation symmetry}{\pw}{prop:sym}
Let $\ell=\alpha\tau+\beta\xi\ne0$, $\ell(\partial)u=0$ and $P(\partial)u=0$, $P$ nonconstant. (i) If $P\not\sim\ell$, $\Ann(u)\ne(P)$. (ii) If $\ell\nmid P$, $u\in\EP$. (iii) Every dictionary containing the monomials of $P$ and $\ell$ has a Gram kernel of dimension $\ge2$.
\end{dimcard}
\begin{proof}
\step{1}{(i)}
$\ell\in\Ann(u)$; if $\Ann(u)=(P)$ then $P\mid\ell$, so $P\sim\ell$.
\step{2}{(ii)}
$\ell$ is irreducible and does not divide $P$, so they are coprime; $\Var(\Ann u)\subseteq\Var(P,\ell)$ is finite and $u\in\EP$ (A2).
\step{3}{(iii)}
$P$ and $\ell$ are independent elements of the annihilator.
\whyworks{a constant-coefficient symmetry is itself an equation, and two independent equations leave only finitely many exponentials.}
\end{proof}
\uses{A1, A2.}

%---------------------------------------------------------------------
\begin{dimcard}{B15}{C3-P4}{Spectral criterion under uniqueness of the expansion}{\pw}{}
Let $P$ be squarefree, $u=\sum_{z\in Z}c_ze^{\langle z,\cdot\rangle}$ in $\mathcal D'(\Omega)$, $Z\subset\Var(P)$ countable, $c_z\ne0$, and assume uniqueness of such expansions (H-uniq). Then $\Ann(u)=\{Q:\ Q|_Z=0\}$, and $\Ann(u)=(P)$ iff $Z\cap\Var(P_i)$ is infinite for every $i$.
\end{dimcard}
\begin{proof}
\step{1}{Annihilators}
$Q(\partial)u=\sum c_zQ(z)e^{\langle z,\cdot\rangle}$, zero iff $Q|_Z=0$ by (H-uniq).
\step{2}{($\Leftarrow$)}
If $Q|_Z=0$ and $Z\cap\Var(P_i)$ is infinite, then $Q$ and $P_i$ have infinitely many common zeros, so $P_i\mid Q$ (A1(V)); hence $P\mid Q$.
\step{3}{($\Rightarrow$)}
If $F=Z\cap\Var(P_i)$ is finite, $Q=(P/P_i)\prod_{z\in F}\ell_z$ with linear $\ell_z(z)=0$ vanishes on $Z$ and is not a multiple of $P$.
\whyworks{with a unique expansion, the exponents are the points where annihilators must vanish.}
\end{proof}
\uses{A1.}

%---------------------------------------------------------------------
\begin{dimcard}{B16}{C3-P5}{Local data}{\pw{} (a) / \cf{} (b)}{prop:local}
(a) For $P=\tau+\xi$, $F(s)=e^s+\varphi(s)$ with a nonzero bump $\varphi\in C_c^\infty(2,3)$ and $u=F(x-t)$: on $(0,1)\times(-5,5)$, $\Ann(u)=(\tau+\xi)$; on $(0,1)^2$, $u=e^{x-t}$ and $\Ann$ is zero-dimensional. (b) For elliptic $P$, $\Ann(u|_{\Omega'})=\Ann(u)$ for every connected $\Omega'\subset\Omega$ (Petrowsky analyticity) \cf.
\end{dimcard}
\begin{proof}
\step{1}{(a)}
If $u\in\EP$ on the large rectangle, $F$ would be an exponential-polynomial on $(-6,5)$, analytic, equal to $e^s$ on $(-1,1)$ hence everywhere, but $F\ne e^s$ on $(2,3)$. On the small square $x-t\in(-1,1)$ and $u=e^{x-t}$.
\step{2}{(b)}
$Q(\partial)u$ solves the elliptic equation, so it is analytic; if it vanishes on $\Omega'$ it vanishes on $\Omega$.
\whyworks{a non-analytic bump can hide outside a smaller window; analyticity forbids it.}
\end{proof}
\uses{A2.}

%=====================================================================
\section{Decidability and variable coefficients}

\begin{dimcard}{B17}{P4-C1}{Restriction does not identify in the Weyl algebra}{\pw}{prop:weyl}
(a) For every nonzero smooth $u$ of one variable, the action $A_1\to\mathrm{End}(A_1\cdot u)$ is injective. (b) For $u=e^x$ and $\mathcal L=\spn\{1,\partial\}$ the restriction to $\mathcal L$ is injective, yet $\partial-1\in\Ann(u)\cap\mathcal L$. (c) The correct object is the evaluation $L\mapsto Lu$, whose kernel is $\Ann(u)\cap\mathcal L$.
\end{dimcard}
\begin{proof}
\step{1}{(a)}
The kernel of the action is a two-sided ideal not containing $1$; $A_1$ is simple, so it is $0$.
\step{2}{(b)}
$(\partial-1)e^x=0$. If $c_0+c_1\partial$ acts as zero on $\C[x]e^x$, applying it to $xe^x$ gives $c_0x+c_1(1+x)=0$, so $c_0=c_1=0$.
\whyworks{in a noncommutative ring, killing the generator does not kill the cyclic module.}
\end{proof}
\uses{none.}

%---------------------------------------------------------------------
\begin{dimcard}{B18}{OP-ALG-T1}{Decidability in the holonomic model}{\modulo{(AC)}}{thm:holo}
\phantomsection\label{dim:P4-T6}%
In the analytic holonomic model M3 (and for analytic solutions with D-finite Cauchy data, \led{P4-T6}), under (AC): ``$L_*u=0$?'' is decidable for every operator with polynomial coefficients; ``$Q(\partial)u\in\EP$?'' is decidable; ``$\Ann(u)=(P)$?'' and ``$\Var(\Ann u)=\Var(P)$?'' are decidable and $G_u$ is computable; a bound $B$ and generators of $\Ann(u)$ are computable.
\end{dimcard}
\begin{proof}
\step{1}{Zero test}
Holonomic functions are closed under the operators involved; a holonomic function given with enough initial values is zero iff those values are zero (OP-ALG-P1), which is a finite check in a computable field under (AC).
\step{2}{EP test}
For $w=Q(\partial)u$ compute eliminating operators in $t$ alone and in $x$ alone; $w\in\EP$ iff their constant-coefficient parts $a^*(\partial_t)$, $b^*(\partial_x)$ kill $w$ (OP-ALG-L4, P2), two zero tests.
\step{3}{Identification and $G_u$}
A3 and A4 with at most $\sum_ie_i$ EP tests.
\step{4}{Bound and generators}
$w_0=G_u(\partial)u\in\EP$ has $\ring u\cap\EP=\ring w_0$ (A5); its eliminants give a finite-dimensional stable space containing $\ring w_0$, hence $B$; A6(ii) and A8(i) give generators.
\whyworks{holonomic functions carry finite certificates of their own equations, which supply the bound that a black box lacks.}
\end{proof}
\uses{A3, A4, A5, A6, A8; research file \texttt{work/OP-ALG}.}

%---------------------------------------------------------------------
\begin{dimcard}{B19}{OP-ALG-T3}{The minimal sufficient datum}{\pw}{}
The minimal datum that makes ``$\Ann(u)=(P)$'' decidable in M2$^+$ is equivalent to a degree $D$ such that, if identification fails, a witness of degree $\le D$ exists. The bound $B$ of A8 is equivalent to knowing $G_u$, and it is strictly stronger than the minimal datum on some regions when $\sum_ie_i\ge2$.
\end{dimcard}
\begin{proof}
\step{1}{$D$ suffices and is computable from any sufficient datum}
Given $D$, compare annihilators in degree $\max(D,\deg P)$ (A8(i)). Conversely, from the answer, run the semi-decision A8(ii), which returns a witness degree when identification fails.
\step{2}{$B\equiv G_u$}
A8(iii) computes $G_u$ from $B$; conversely, from $G_u$, $w_0=G_u(\partial)u$ is an exponential-polynomial and searching annihilators $a(\tau)$, $b(\xi)$ of $w_0$ gives $B=\deg a\cdot\deg b$.
\whyworks{deciding needs only the degree of a witness, while $B$ describes the whole exponential part.}
\end{proof}
\uses{A5, A8.}

%---------------------------------------------------------------------
\begin{dimcard}{B20}{OP-ALG-T4}{Saturated identification criterion}{\pw}{thm:sat}
Let $\mathfrak D$ be a ring of operators $\mathfrak A[\partial_t]$ with $\mathfrak A$ a domain of operators in $x$, $P$ monic in $\partial_t$ of degree $m$, $Pu=0$. Equivalent: (a) $\Ann_{\mathfrak D}(u)=\mathfrak DP$; (b) the $t$-jet is free over $\mathfrak A$; (c) $\mathrm{rank}_{\mathfrak A}(\mathfrak D u)=m$; (d) no $Q\notin\mathfrak DP$ has $Qu$ in the $\mathfrak A$-torsion.
\end{dimcard}
\begin{proof}
\step{1}{Division}
$Q=AP+Q_0$ with $\deg_{\partial_t}Q_0<m$, $Qu=Q_0u$, and $Q\in\mathfrak DP$ iff $Q_0=0$.
\step{2}{The equivalences}
(a)$\iff$(b): a strictly larger annihilator is a nonzero $Q_0$ of degree $<m$ killing $u$. (b)$\iff$(c): $\mathfrak Du\cong\mathfrak A^m/N$ with $N$ the module of relations, torsion-free. (b)$\iff$(d): a relation is a $Q_0$ with $Q_0u=0$; conversely $bQ_0u=0$ with $b\ne0$ gives the relation $bQ_0$.
\whyworks{as in B9, modulo $P$ everything is a combination of the first $m$ time derivatives.}
\end{proof}
\uses{B9.}

%---------------------------------------------------------------------
\begin{dimcard}{B21}{OP-ALG-C1}{With variable coefficients the finite test fails}{\pw}{}
With $\theta=x\partial_x$, $P_1=\partial_t-\theta$, $P_2=\partial_t-\theta+1$, $P=P_1P_2$, $u_2=e^{-t}\exp(x^2e^{2t})$ and $u^{(\lambda)}=(1+\lambda\partial_x)u_2$: $\Ann(u^{(\lambda)})\ne\mathfrak DP$ for every $\lambda$, yet neither partial quotient $P_1u^{(\lambda)}=-u_2$ nor $P_2u^{(\lambda)}=\lambda\partial_xu_2$ is torsion. So the literal transfer of A3(a) is false. \textbf{(OP-ALG-P5)} In the full Weyl algebra $A_2$, a holonomic solution never identifies its equation.
\end{dimcard}
\begin{proof}
\step{1}{The identities}
$P_1P_2=P_2P_1$, $P_2-P_1=1$, $P_1\partial_x=\partial_xP_2$, so $(1+\lambda\partial_x)P=P_1Q_\lambda$ with $Q_\lambda=P_2+\lambda\partial_xP_1$, and $Q_\lambda u^{(\lambda)}=0$ with $Q_\lambda\notin\mathfrak DP$.
\step{2}{(P5)}
A holonomic $u$ is killed by a nonzero $L$ in $x$ alone, so $u$ is torsion and B20(d) fails with $Q=1$.
\whyworks{without a gcd, annihilators of noncommutative operators are not controlled by finitely many quotients.}
\end{proof}
\uses{B20.}

\section{Results used without proof}

\listcard{C3-P1[gen]}{Malgrange density: on convex $\Omega$ exponential solutions are dense in $\ker P$; finitely many continuous functionals never certify identification.}{\cf}{Malgrange, Ann.\ Inst.\ Fourier 6, Chap.~I \S2 Thm~2 (primary source checked).}
\listcard{II-T9}{With variable coefficients in $\C(x)[\partial]$, identification needs the pair (module, cyclic vector).}{\cf}{Ore 1933; van der Put--Singer.}
\listcard{R6-T3}{``The companion matrix is normal iff $p=s^m-c$, $|c|=1$'' for every $m$.}{\falsest}{Counterexample $s^2+s-1$ in B8; true for $m\ge3$.}
